\begin{lemma}[Composition of quadratic envelopes]
Let $F$ be a finite-dimensional real inner-product space, let $I\subseteq\R$ be
open, let $\Omega=\{(x,t):t\in I\}$, and let $f$ be a real function on space-time
which is continuous on $\Omega$ and has subquadratic spatial growth on compact
time intervals of $I$ in the sense of \noderef{SubquadraticGrowth}. Then,
whenever $\ell,h>0$ and the time arguments $t$, $t+h$, $t+\ell+h$ all lie in
$I$,
\begin{equation}\label{eq:envelope-composition}
  \Penv_\ell\bigl(\Penv_hf\bigr)(x,t)=\Penv_{\ell+h}f(x,t),
\end{equation}
and likewise, whenever $\ell,h>0$ and $t$, $t-h$, $t-\ell-h$ all lie in $I$,
\[
  \Menv_\ell\bigl(\Menv_hf\bigr)(x,t)=\Menv_{\ell+h}f(x,t),
\]
with the envelopes of \noderef{QuadEnvP} and \noderef{QuadEnvM}. Here
$\Penv_hf$ and $\Menv_hf$ are regarded as functions of their center and time,
the time being shifted by $h$ in the respective direction.
\end{lemma}
\begin{proof}
Throughout, $z=(x,t)$ is a space-time point with $x\in F$ and $t\in\R$, and
$\ell,h>0$ are fixed. For a fixed pair $(x,w)\in F\times F$ we write
\[
  Q(y)=\frac{|y-x|^2}{2\ell}+\frac{|w-y|^2}{2h}\qquad(y\in F)
\]
for the combined cost. By the convention of the preamble a supremum over a set
which is empty or unbounded above, and an infimum over a set which is empty or
unbounded below, take the junk value $0$; accordingly every extremum occurring
below is shown to be over a nonempty set carrying the relevant one-sided bound
\emph{before} that extremum is used. Steps 1--5 prove
\eqref{eq:envelope-composition} and Step 6 proves the backward identity; Step 6
also explains why the backward identity is not the forward one applied to $-f$.

\emph{Step 1: the parallel-axis identity.}
Fix $x,w\in F$ and set
\[
  a=\frac1{2\ell},\qquad b=\frac1{2h},\qquad
  y_*=\frac{a\,x+b\,w}{a+b} .
\]
Since $a+b=\dfrac{\ell+h}{2\ell h}$, we have $\dfrac{a}{a+b}=\dfrac{h}{\ell+h}$
and $\dfrac{b}{a+b}=\dfrac{\ell}{\ell+h}$, so
$y_*=\dfrac{h\,x+\ell\,w}{\ell+h}$ is the convex combination of $x$ and $w$
with weights $h/(\ell+h)$ and $\ell/(\ell+h)$. Now $Q(y)=a|y-x|^2+b|y-w|^2$,
and expanding both squares by bilinearity and symmetry of the inner product,
\[
  Q(y)=(a+b)|y|^2-2\langle y,\,a\,x+b\,w\rangle+a|x|^2+b|w|^2 .
\]
Because $a+b>0$ we may complete the square, obtaining
\[
  Q(y)=(a+b)\Bigl|y-\frac{a\,x+b\,w}{a+b}\Bigr|^2
    -\frac{|a\,x+b\,w|^2}{a+b}+a|x|^2+b|w|^2 .
\]
For the constant term, $|a\,x+b\,w|^2=a^2|x|^2+2ab\langle x,w\rangle+b^2|w|^2$,
so
\[
  (a+b)\bigl(a|x|^2+b|w|^2\bigr)-|a\,x+b\,w|^2
   =ab|x|^2-2ab\langle x,w\rangle+ab|w|^2
   =ab\,|w-x|^2 ,
\]
the $a^2|x|^2$ and $b^2|w|^2$ terms cancelling, and hence
\[
  a|x|^2+b|w|^2-\frac{|a\,x+b\,w|^2}{a+b}=\frac{ab}{a+b}\,|w-x|^2 .
\]
Finally $\dfrac{ab}{a+b}=\dfrac{1}{4\ell h}\cdot\dfrac{2\ell h}{\ell+h}
=\dfrac1{2(\ell+h)}$. Substituting the last two displays and
$a+b=\frac{\ell+h}{2\ell h}$ into the completed square gives, for every
$y\in F$,
\begin{equation}\label{eq:parallel-axis}
  \frac{|y-x|^2}{2\ell}+\frac{|w-y|^2}{2h}
   =\frac{|w-x|^2}{2(\ell+h)}+\frac{\ell+h}{2\ell h}\,|y-y_*|^2 .
\end{equation}
Three consequences are used below, each immediate from
\eqref{eq:parallel-axis}.

(1a) Since $\frac{\ell+h}{2\ell h}>0$ and $|y-y_*|^2\ge0$, for \emph{every}
$y\in F$
\[
  \frac{|y-x|^2}{2\ell}+\frac{|w-y|^2}{2h}\ \ge\ \frac{|w-x|^2}{2(\ell+h)} .
\]

(1b) Taking $y=y_*$ in \eqref{eq:parallel-axis} gives
$Q(y_*)=\frac{|w-x|^2}{2(\ell+h)}$. Hence $\{Q(y):y\in F\}$ is nonempty and,
by (1a), bounded below by $\frac{|w-x|^2}{2(\ell+h)}$, which is one of its
elements; therefore
\[
  \inf_{y\in F}Q(y)=\frac{|w-x|^2}{2(\ell+h)} ,
\]
the infimum being attained at $y_*$ and, since $\frac{\ell+h}{2\ell h}>0$
forces $|y-y_*|^2=0$ for any other minimizer, only there.

(1c) Consistently with (1b), $y_*-x=\frac{\ell}{\ell+h}(w-x)$ and
$w-y_*=\frac{h}{\ell+h}(w-x)$ by direct computation from
$y_*=\frac{h x+\ell w}{\ell+h}$, and substituting these two into $Q$ gives
$\frac1{2\ell}\frac{\ell^2|w-x|^2}{(\ell+h)^2}
 +\frac1{2h}\frac{h^2|w-x|^2}{(\ell+h)^2}
 =\frac{|w-x|^2}{2(\ell+h)^2}(\ell+h)=\frac{|w-x|^2}{2(\ell+h)}$.

\emph{Step 2: the single-cost family is nonempty and bounded above.}
Assume from here to the end of Step 5 that $t$, $t+h$ and $t+\ell+h$ lie in
$I$. Put $s=t+\ell+h$ and define
\[
  G(w)=f(w,s)-\frac{|w-x|^2}{2(\ell+h)}\qquad(w\in F).
\]
We claim there is $B\in\R$ with $G(w)\le B$ for all $w\in F$. The singleton
$J=\{s\}$ is compact and $J\subseteq I$ because $s=t+\ell+h\in I$ by
hypothesis, so \noderef{SubquadraticGrowth} applies to $J$ with
$\varepsilon=\frac1{8(\ell+h)}>0$ and yields $M>0$ such that
\[
  |f(w,s)|\le\frac{|w|^2}{8(\ell+h)}\qquad(|w|\ge M).
\]
Set $R=\max\{M,2|x|\}$. If $|w|\ge R$ then $|x|\le|w|/2$, so by the triangle
inequality $|w-x|\ge|w|-|x|\ge|w|/2\ge0$ and hence $|w-x|^2\ge|w|^2/4$;
combining,
\[
  \frac{|w-x|^2}{2(\ell+h)}\ \ge\ \frac{|w|^2}{8(\ell+h)}\ \ge\ f(w,s),
\]
so $G(w)\le0$. If $|w|\le R$, note that $F$ is finite dimensional and
$\overline B_R=\{w\in F:|w|\le R\}$ is closed and bounded, hence compact, and
nonempty because $R\ge M>0$ gives $0\in\overline B_R$; the map
$w\mapsto(w,s)$ is continuous and takes $\overline B_R$ into
$\Omega=\{(y,\sigma):\sigma\in I\}$ because $s\in I$, so $w\mapsto f(w,s)$ is
continuous on $\overline B_R$ by the continuity hypothesis on $f$, and
$w\mapsto\frac{|w-x|^2}{2(\ell+h)}$ is continuous. Thus $G$ is continuous on
the nonempty compact set $\overline B_R$ and attains a finite maximum $B_0$
there. With
$B=\max\{B_0,0\}$ we get $G(w)\le B$ for every $w\in F$. The family is also
nonempty, $F$ containing $x$.

\emph{Step 3: every extremum in the iterated envelope is legitimate.}
Define
\[
  H(y,w)=f(w,s)-\frac{|w-y|^2}{2h}-\frac{|y-x|^2}{2\ell}
  \qquad(y,w\in F).
\]
Applying (1a) at the pair $(x,w)$, we have for every $y\in F$
\[
  \frac{|y-x|^2}{2\ell}+\frac{|w-y|^2}{2h}\ \ge\ \frac{|w-x|^2}{2(\ell+h)} ,
\]
so for all $y,w\in F$
\begin{equation}\label{eq:two-cost-dominated}
  H(y,w)\ \le\ f(w,s)-\frac{|w-x|^2}{2(\ell+h)}=G(w)\ \le\ B .
\end{equation}
Consequently:
\begin{itemize}
\item for each fixed $y\in F$ the set
$\bigl\{f(w,s)-\frac{|w-y|^2}{2h}:w\in F\bigr\}$ is nonempty (take $w=y$) and
bounded above by $B+\frac{|y-x|^2}{2\ell}$, by \eqref{eq:two-cost-dominated};
\item the set $\{H(y,w):(y,w)\in F\times F\}$ is nonempty (take $y=w=x$) and
bounded above by $B$;
\item for each fixed $w\in F$ the set $\{H(y,w):y\in F\}$ is nonempty (take
$y=x$) and bounded above by $B$;
\item the set $\bigl\{f(w,s)-\frac{|w-x|^2}{2(\ell+h)}:w\in F\bigr\}$ is
nonempty and bounded above by $B$, by Step 2.
\end{itemize}
By \noderef{QuadEnvP} with length $h\ne0$, at the center $y$ and time
$t+\ell$, the envelope $\Penv_hf(y,t+\ell)$ is the supremum of the first of
these sets, since the time argument of $f$ is $(t+\ell)+h=s$; the first bullet
makes it a genuine supremum, in particular finite. Note that $f$ is evaluated
only at time $s=t+\ell+h$, so no information about $t+\ell$ is required here or
anywhere below. By \noderef{QuadEnvP} with length $\ell\ne0$ at the center
$(x,t)$,
\begin{equation}\label{eq:outer-envelope}
  \Penv_\ell\bigl(\Penv_hf\bigr)(x,t)
   =\sup_{y\in F}\left\{\Penv_hf(y,t+\ell)-\frac{|y-x|^2}{2\ell}\right\},
\end{equation}
and this outer set is nonempty and bounded above by $B$: for each $y$,
subtracting the constant $\frac{|y-x|^2}{2\ell}$ from the first bullet's upper
bound gives $\Penv_hf(y,t+\ell)-\frac{|y-x|^2}{2\ell}\le B$.

\emph{Step 4: the interchange, in three separate moves.}

\emph{(4i) Translating the inner supremum by a $w$-independent constant.} Fix
$y\in F$. For a nonempty set $S\subseteq\R$ bounded above and $c\in\R$ one has
$\sup\{u-c:u\in S\}=\sup S-c$, since $u\le\sup S$ for all $u\in S$ gives
$\le$, and $\sup S\le\sup\{u-c:u\in S\}+c$ follows because
$\sup\{u-c:u\in S\}+c$ is an upper bound for $S$. Applying this with $S$ the
first bullet's set and $c=\frac{|y-x|^2}{2\ell}$,
\[
  \Penv_hf(y,t+\ell)-\frac{|y-x|^2}{2\ell}
   =\sup_{w\in F}\left\{f(w,s)-\frac{|w-y|^2}{2h}
     -\frac{|y-x|^2}{2\ell}\right\}
   =\sup_{w\in F}H(y,w).
\]

\emph{(4ii) Iterated supremum equals supremum over the product.} All three
families in the second and third bullets are nonempty and bounded above, so
$S^*=\sup_{(y,w)\in F\times F}H(y,w)$, $\sup_{w}H(y,w)$ for each $y$, and
$\sup_{y}H(y,w)$ for each $w$ all exist in $\R$. For each $y$ and each $w$,
$H(y,w)\le S^*$, so $\sup_wH(y,w)\le S^*$; the set
$\{\sup_wH(y,w):y\in F\}$ is therefore nonempty and bounded above by $S^*$,
giving $\sup_y\sup_wH(y,w)\le S^*$. Conversely, for any $(y,w)$,
$H(y,w)\le\sup_{w'}H(y,w')\le\sup_{y'}\sup_{w'}H(y',w')$, so
$\sup_{y'}\sup_{w'}H(y',w')$ is an upper bound for the product family and
$S^*\le\sup_y\sup_wH(y,w)$. Hence
$\sup_y\sup_wH(y,w)=S^*$, and the identical argument with the roles of the two
coordinates exchanged gives $\sup_w\sup_yH(y,w)=S^*$. Combining with
\eqref{eq:outer-envelope} and (4i),
\[
  \Penv_\ell\bigl(\Penv_hf\bigr)(x,t)
   =\sup_{y\in F}\sup_{w\in F}H(y,w)=S^*=\sup_{w\in F}\sup_{y\in F}H(y,w).
\]

\emph{(4iii) Pulling out the $y$-independent term and passing to the
infimum.} Fix $w\in F$. The quantity $f(w,s)$ does not depend on $y$, and
$H(y,w)=f(w,s)-Q(y)$ with $Q$ as in Step 1. By (1b) the set $\{Q(y):y\in F\}$
is nonempty and bounded below, so for a nonempty $T\subseteq\R$ bounded below
we may use $\sup\{-v:v\in T\}=-\inf T$: indeed $-v\le-\inf T$ for every
$v\in T$ gives $\le$, and any upper bound $\beta$ for $\{-v:v\in T\}$ makes
$-\beta$ a lower bound for $T$, so $-\beta\le\inf T$, i.e.
$-\inf T\le\beta$; as $\sup\{-v:v\in T\}$ is the least upper bound, the two
halves give equality. Applying this with $T=\{Q(y):y\in F\}$, which is nonempty
and bounded below by (1b),
\[
  \sup_{y\in F}\bigl(-Q(y)\bigr)=-\inf_{y\in F}Q(y)
   =-\frac{|w-x|^2}{2(\ell+h)} ,
\]
the second equality by (1b). In particular $\{-Q(y):y\in F\}$ is nonempty and
bounded above, so the translation rule of (4i) applies to it with
$c=-f(w,s)$; since $H(y,w)=\bigl(-Q(y)\bigr)-\bigl(-f(w,s)\bigr)$, that rule
gives
\[
  \sup_{y\in F}H(y,w)=\sup_{y\in F}\bigl(-Q(y)\bigr)+f(w,s)
   =f(w,s)-\frac{|w-x|^2}{2(\ell+h)}=G(w).
\]

\emph{Step 5: conclusion of the forward identity.}
Combining (4ii) and (4iii),
\[
  \Penv_\ell\bigl(\Penv_hf\bigr)(x,t)
   =\sup_{w\in F}G(w)
   =\sup_{w\in F}\left\{f(w,t+\ell+h)-\frac{|w-x|^2}{2(\ell+h)}\right\},
\]
a supremum over a nonempty set bounded above by Step 2. On the other hand
$\ell+h>0$, so by \noderef{QuadEnvP} with length $\ell+h\ne0$ at the center
$(x,t)$,
\[
  \Penv_{\ell+h}f(x,t)
   =\sup_{w\in F}\left\{f\bigl(w,t+(\ell+h)\bigr)
      -\frac{|w-x|^2}{2(\ell+h)}\right\},
\]
and $t+(\ell+h)=t+\ell+h=s$, so the two right-hand sides are suprema of the
same set. This is \eqref{eq:envelope-composition}. (The argument used only
$t+\ell+h\in I$; the hypotheses $t\in I$ and $t+h\in I$ belong to the
statement's admissibility convention and are not needed for this identity.)

\emph{Step 6: the backward identity, written out.}
Assume now $\ell,h>0$ and $t$, $t-h$, $t-\ell-h\in I$, and put
$\sigma=t-\ell-h$. This is \emph{not} the forward case applied to $-f$:
unwinding \noderef{QuadEnvP} for $-f$ gives
$-\Penv_\ell(-f)(x,t)=\inf_{y}\{f(y,t+\ell)+\frac{|y-x|^2}{2\ell}\}$, whose
time argument is $t+\ell$, whereas \noderef{QuadEnvM} shifts time by $-\ell$;
reducing one to the other would additionally require reflecting the time
variable about a base point, which replaces $I$ by its reflection and so is
unavailable under hypotheses stated for the fixed set $I$. Note also that the
cost signs are \emph{not} reversed below: the two cost terms are added, exactly
as in $Q$, and the same infimum (1b) is used, with no appeal to
$\sup(-T)=-\inf T$.

By \noderef{QuadEnvM} with length $h\ne0$ at the center $y$ and time
$t-\ell$, the time argument of $f$ is $(t-\ell)-h=\sigma$, so
\[
  \Menv_hf(y,t-\ell)
   =\inf_{w\in F}\left\{f(w,\sigma)+\frac{|w-y|^2}{2h}\right\},
\]
and by \noderef{QuadEnvM} with length $\ell\ne0$ at $(x,t)$,
\[
  \Menv_\ell\bigl(\Menv_hf\bigr)(x,t)
   =\inf_{y\in F}\left\{\Menv_hf(y,t-\ell)+\frac{|y-x|^2}{2\ell}\right\}.
\]
Set $\widetilde G(w)=f(w,\sigma)+\frac{|w-x|^2}{2(\ell+h)}$ and
$\widetilde H(y,w)=f(w,\sigma)+\frac{|w-y|^2}{2h}+\frac{|y-x|^2}{2\ell}$.
Because an infimum over a set unbounded below is likewise the junk value $0$,
this direction needs its own \emph{lower} bound, which is the mirror of Step 2.
The singleton $\{\sigma\}$ is compact and contained in $I$ since
$\sigma=t-\ell-h\in I$ by hypothesis, so \noderef{SubquadraticGrowth} with
$\varepsilon=\frac1{8(\ell+h)}>0$ gives $M>0$ with
$|f(w,\sigma)|\le\frac{|w|^2}{8(\ell+h)}$ for $|w|\ge M$. With
$R=\max\{M,2|x|\}$, if $|w|\ge R$ then as in Step 2
$|w-x|^2\ge|w|^2/4$, hence
\[
  \frac{|w-x|^2}{2(\ell+h)}\ \ge\ \frac{|w|^2}{8(\ell+h)}\ \ge\ -f(w,\sigma),
\]
so $\widetilde G(w)\ge0$; and on $\overline B_R=\{w\in F:|w|\le R\}$, which is
closed, bounded and nonempty exactly as in Step 2, hence compact and nonempty,
the function $\widetilde G$ is continuous, because $\sigma\in I$ puts
$(w,\sigma)$ in $\Omega$ for every such $w$, so $\widetilde G$ attains a finite
minimum $b_0$ there. With
$b=\min\{b_0,0\}$ we get $\widetilde G(w)\ge b$ for all $w\in F$. By (1a),
$\frac{|y-x|^2}{2\ell}+\frac{|w-y|^2}{2h}\ge\frac{|w-x|^2}{2(\ell+h)}$, so
\[
  \widetilde H(y,w)\ \ge\ \widetilde G(w)\ \ge\ b\qquad(y,w\in F).
\]
Hence, mirroring Step 3: for each fixed $y$ the set
$\{f(w,\sigma)+\frac{|w-y|^2}{2h}:w\in F\}$ is nonempty ($w=y$) and bounded
below by $b-\frac{|y-x|^2}{2\ell}$, so $\Menv_hf(y,t-\ell)$ is a genuine
infimum; the families $\{\widetilde H(y,w):(y,w)\in F\times F\}$,
$\{\widetilde H(y,w):y\in F\}$ for fixed $w$, and
$\{\widetilde G(w):w\in F\}$ are nonempty and bounded below by $b$; and the
outer set $\{\Menv_hf(y,t-\ell)+\frac{|y-x|^2}{2\ell}:y\in F\}$ is nonempty
and bounded below by $b$.

The three moves of Step 4 now run with every supremum replaced by an infimum
and each inequality reversed, the corresponding elementary facts being
$\inf\{u+c:u\in S\}=\inf S+c$ for nonempty $S$ bounded below and $c\in\R$, and
the equality of the iterated infimum with the infimum over the product for a
nonempty family bounded below (same two-inequality argument as in (4ii), with
lower bounds in place of upper bounds). This gives
\[
  \Menv_\ell\bigl(\Menv_hf\bigr)(x,t)
   =\inf_{y\in F}\inf_{w\in F}\widetilde H(y,w)
   =\inf_{w\in F}\inf_{y\in F}\widetilde H(y,w),
\]
and for fixed $w$, since $f(w,\sigma)$ does not depend on $y$ and
$\widetilde H(y,w)=f(w,\sigma)+Q(y)$,
\[
  \inf_{y\in F}\widetilde H(y,w)=f(w,\sigma)+\inf_{y\in F}Q(y)
   =f(w,\sigma)+\frac{|w-x|^2}{2(\ell+h)}=\widetilde G(w),
\]
by the translation rule and (1b) -- the same infimum as in the forward case,
with no sign change. Therefore
\[
  \Menv_\ell\bigl(\Menv_hf\bigr)(x,t)
   =\inf_{w\in F}\left\{f(w,t-\ell-h)+\frac{|w-x|^2}{2(\ell+h)}\right\},
\]
an infimum over a nonempty set bounded below. Finally $\ell+h>0$, so by
\noderef{QuadEnvM} with length $\ell+h\ne0$ at $(x,t)$,
\[
  \Menv_{\ell+h}f(x,t)
   =\inf_{w\in F}\left\{f\bigl(w,t-(\ell+h)\bigr)
      +\frac{|w-x|^2}{2(\ell+h)}\right\},
\]
and $t-(\ell+h)=t-\ell-h=\sigma$, so the two infima are over the same set and
$\Menv_\ell(\Menv_hf)(x,t)=\Menv_{\ell+h}f(x,t)$.
\end{proof}
